%% TRANSCRIBED (PDF 98-110 = printed III-7..III-19)
%% NOTE: printed pages III-6 and III-7 are BLANK placeholder pages in the
%% working PDF (Taylor & Francis ebook defect; local PDF 97-98 contain only
%% the publisher logo).  The beginning of this section on III-7 (the opening
%% paragraph of 2.1, up to ``be a place'') was recovered from the textually
%% identical 1989 Addison-Wesley photo-reprint of the 1968 Benjamin edition
%% (public McGill scan) and was checked to continue verbatim into III-8 of
%% the working copy.
%% TODO-VERIFY(p098): the III-7 passage below follows the 1989 reprint scan,
%% not the working PDF (whose page 98 is blank).

\sechead{2}{The global case}{THE GLOBAL CASE}

\subsechead{2.1}{Definitions}

We now go back to the notations of chap.~\ref{chap:II}, i.e. $K$ denotes a
number field.  Let $\ell$ be a prime number and let
\[
\rho\colon \Gal(\bar K/K)^{\mathrm{ab}} \to \Aut(V_\ell)
\]
be an abelian \textit{l}-adic representation of $K$.  Let $v\in\Sigma_K$
be a place of $K$ of residue characteristic $\ell$ and let
$D_v\subset \Gal(\bar K/K)^{\mathrm{ab}}$ be the corresponding
decomposition group.  This group is a quotient of the local Galois group
$\Gal(\bar K_v/K_v)^{\mathrm{ab}}$ (these two groups are, in fact,
isomorphic, but we do not need this here).  Hence, we get by composition
an \textit{l}-adic representation of $K_v$
\[
\rho_v\colon \Gal(\bar K_v/K_v)^{\mathrm{ab}} \to D_v
\xrightarrow{\ \rho\ } \Aut(V_\ell)\,.
\]

\begin{dfn}
The representation $\rho$ is said to be locally algebraic if all the
local representations $\rho_v$, with $p_v=\ell$, are locally algebraic
(in the sense defined in \S\ref{sec:III.1.1}, with $p=\ell$).
\end{dfn}

It is convenient to reformulate this definition, using the torus
$T=\mathrm{R}_{K/\Q}(\mathbf{G}_{m/K})$ of Chap.~\ref{chap:II},
\S\ref{sec:II.1.1}.  Let $T/\Q_\ell = T\times_\Q \Q_\ell$ be the
$\Q_\ell$-torus obtained from $T$ by extending the ground field from $\Q$
to $\Q_\ell$.  We have
\[
T/\Q_\ell(\Q_\ell) = (K\tens\Q_\ell)^* = K_\ell^*\,,
\]
where $K_\ell = K\tens\Q_\ell$.

Let $I$ be the idèle group of $K$, cf.~Chap.~\ref{chap:II},
\S\ref{sec:II.2.1}.  The injection $K_\ell^*\to I$, followed by the class
field homomorphism $i\colon I\to\Gal(\bar K/K)^{\mathrm{ab}}$, defines a
homomorphism
\[
i_\ell\colon K_\ell^* \to \Gal(\bar K/K)^{\mathrm{ab}}\,.
\]

\begin{proposition}
The representation $\rho$ is locally algebraic if and only if there
exists an algebraic morphism
\[
f\colon T/\Q_\ell \to \GL_{V_\ell}
\]
such that $\rho\circ i_\ell(x) = f(x^{-1})$ for all $x\in K_\ell^*$
close enough to $1$.
\end{proposition}

(Note that, as in the local case, the above condition determines $f$
uniquely; one says it is the algebraic morphism \emph{associated with}
$\rho$.)

Since $K\tens_\Q \Q_\ell = \prod_{v|\ell} K_v$, we have
\[
T/\Q_\ell = \prod_{v|\ell} T_v\,,
\]
where $T_v$ is the $\Q_\ell$-torus defined by $K_v$,
cf.~\S\ref{sec:III.1.1}.  The proposition follows from this decomposition.

\runhead{Exercise}

Give a criterion for local algebraicity analogous to the one of
Prop.~\ref{prop:III.1.1.2} of \S\ref{sec:III.1.1}.

\subsechead{2.2}{Modulus of a locally algebraic abelian representation}

Let $\rho\colon \Gal(\bar K/K)^{\mathrm{ab}}\to\Aut(V_\ell)$ be as above;
by composition with the class field homomorphism
$i\colon I\to\Gal(\bar K/K)^{\mathrm{ab}}$, $\rho$ defines a homomorphism
$\rho\circ i\colon I\to\Aut(V_\ell)$.

We assume that $\rho$ is locally algebraic and we denote by $f$ the
associated algebraic morphism $T/\Q_\ell\to\GL_{V_\ell}$.

\begin{dfn}
Let $m$ be a modulus (chap.~\ref{chap:II}, \S\ref{sec:II.1.1}).  One says
that $\rho$ is defined mod $m$ (or that $m$ is a modulus of definition
for $\rho$) if
\begin{rlist}
\item $\rho\circ i$ is trivial on $U_{v,m}$ when $p_v\neq\ell$.
\item $\rho\circ i_\ell(x) = f(x^{-1})$ for
$x\in\prod_{v|\ell}U_{v,m}$
\end{rlist}
(Note that $\prod_{v|\ell}U_{v,m}$ is an open subgroup of
$K_\ell^* = T/\Q_\ell(\Q_\ell)$.)
\end{dfn}

In order to prove the existence of a modulus of definition, we need the
following auxiliary result:

\begin{proposition}
%% The book prints this PROPOSITION without a number, but refers to it as
%% ``Prop. 1'' in the proof of Proposition 2 below; those references are
%% therefore typed literally (a \ref would print the hidden statement
%% counter, not ``1'').
Let $H$ be a Lie group over $\Q_\ell$ (resp. $\R$) and let $\alpha$ be a
continuous homomorphism of the idèle group $I$ into $H$.
\begin{alist}
\item If $p_v\neq\ell$ (resp. $p_v\neq\infty$), the restriction of
$\alpha$ to $K_v^*$ is equal to $1$ on an open subgroup of $K_v^*$.
\item The restriction of $\alpha$ to the unit group $U_v$ of $K_v^*$ is
equal to $1$ for almost all $v$'s.
\end{alist}
\end{proposition}

Part (a) follows from the fact that $K_v^*$ is a $p_v$-adic Lie group and
that a homomorphism of a $p$-adic Lie group into an \textit{l}-adic one
is locally equal to $1$ if $p\neq\ell$.

To prove (b), let $N$ be a neighborhood of $1$ in $H$ which contains no
finite subgroup except $\{1\}$; the existence of such an $N$ is classical
for real Lie groups, and quite easy to prove for \textit{l}-adic ones.
By definition of the idèle topology, $\alpha(U_v)$ is contained in $N$
for almost all $v$'s.  But (a) shows that, if $p_v\neq\ell$, the group
$\alpha(U_v)$ is finite; hence $\alpha(U_v)=\{1\}$ for almost all $v$'s,
q.e.d.

\begin{corollary}
Any abelian \textit{l}-adic representation of $K$ is unramified outside a
finite set of places.
\end{corollary}

This follows from (b) applied to the homomorphism $\alpha$ of $I$ induced
by the given representation, since the $\alpha(U_v)$ are known to be the
inertia subgroups.

\runhead{Remark}

This does not extend to non-abelian representations (even solvable ones),
cf.~Exercise.

\begin{proposition}[2]\label{prop:III.2.2.2}
Every locally algebraic abelian \textit{l}-adic representation has a
modulus of definition.
\end{proposition}

Let $\rho\colon\Gal(\bar K/K)^{\mathrm{ab}}\to\Aut(V_\ell)$ be the given
representation and $f$ the associated morphism of $T/\Q_\ell$ into
$\GL_{V_\ell}$.  Let $X$ be the set of places $v\in\Sigma_K$, with
$p_v\neq\ell$, for which $\rho$ is ramified; the corollary to Prop.~1
shows that $X$ is finite.  By Prop.~1, (a), we can choose a modulus $m$
such that $\rho\circ i\colon I\to\Aut(V_\ell)$ is trivial on all the
$U_{v,m}$, $v\in X$.  Enlarging $m$ if necessary, we can assume that
$\rho\circ i_\ell(x)=f(x^{-1})$ for $x\in\prod_{p_v=\ell}U_{v,m}$.  Hence,
$m$ is a modulus of definition for $\rho$.

\runhead{Remark}

It is easy to show that there is a smallest modulus of definition for
$\rho$; it is called the \emph{conductor} of $\rho$.

\runhead{Exercise}

Let $z_1,\ldots,z_n,\ldots\in K^*$.  For each $n$, let $E_n$ be the
subfield of $\bar K$ generated by all the $\ell^n$-th roots of the
element $z_1z_2^\ell\ldots z_n^{\ell^{n-1}}$.
\begin{alist}
\item Show that $E_n$ is a Galois extension of $K$, containing the
$\ell^n$-th roots of unity and that its Galois group is isomorphic to a
subgroup of the affine group
$\begin{pmatrix} * & * \\ 0 & 1 \end{pmatrix}$ in
$\GL(2,\Z/\ell^n\Z)$.
\item Let $E$ be the union of the $E_n$'s.  Show that $E$ is a Galois
extension of $K$, whose Galois group is a closed subgroup of the affine
group relative to $\Z_\ell$.
\item Give an example where $E$ (and hence the corresponding
$2$-dimensional \textit{l}-adic representation) is ramified at all places
of $K$.
\end{alist}

\subsechead{2.3}{Back to $S_m$}

Let $m$ be a modulus of $K$ and let
\[
\phi\colon S_m/\Q_\ell \to \GL_{V_\ell}
\]
be a linear representation of $S_m/\Q_\ell$.  Let
\[
\phi_\ell\colon \Gal(\bar K/K)^{\mathrm{ab}} \to \Aut(V_\ell)
\]
be the corresponding \textit{l}-adic representation
(cf.~chap.~\ref{chap:II}, \S\ref{sec:II.2.5}.).

\begin{theorem}[1]\label{thm:III.2.3.1}
The representation $\phi_\ell$ is locally algebraic and defined mod $m$.
The associated algebraic morphism
\[
f\colon T/\Q_\ell \to \GL_{V_\ell}
\]
is $\phi\circ\pi$, where $\pi$ denotes the canonical morphism of $T$ into
$S_m$ (cf.~chap.~\ref{chap:II}, \S\ref{sec:II.2.2}).
\end{theorem}

This is trivial from the construction of $\phi_\ell$ as
$\phi\circ\varepsilon_\ell$ (chap.~\ref{chap:II}, \S\ref{sec:II.2.5}) and
the corresponding properties of $\varepsilon_\ell$
(chap.~\ref{chap:II}, \S\ref{sec:II.2.3}).

The converse of Theorem~\ref{thm:III.2.3.1} is true.  We state it only
for the case of \emph{rational} representations:

\begin{theorem}[2]\label{thm:III.2.3.2}
%% The subscript of V and phi is a small round glyph in the typewriter
%% original; it is read as the digit ``0'' (Q-structure index), cf. the
%% same notation in chap.~II, 2.5.
Let $\rho\colon\Gal(\bar K/K)^{\mathrm{ab}}\to\Aut(V_\ell)$ be an abelian
\textit{l}-adic representation of the number field $K$.  Assume $\rho$ is
rational (chap.~\ref{chap:I}, \S\ref{sec:I.2.3}) and is locally algebraic
with $m$ as a modulus of definition (cf.~\S\ref{sec:III.2.2}).  Then,
there exist a $\Q$-vector subspace $V_0$ of $V_\ell$, with
$V_\ell = V_0\tens_\Q\Q_\ell$, and a morphism
$\phi_0\colon S_m\to\GL_{V_0}$ of $\Q$-algebraic groups such that $\rho$
is equal to the \textit{l}-adic representation $\phi_\ell$ associated to
$\phi_0$ (cf.~chap.~\ref{chap:II}, \S\ref{sec:II.2.5}).
\end{theorem}

(The condition $V_\ell = V_0\tens_\Q\Q_\ell$ means that $V_0$ is a
``$\Q$-structure'' on $V_\ell$, cf.~Bourbaki \emph{Alg.}, chap.~II,
3\textsuperscript{rd} ed.)

\begin{proof}
Let $r\colon T/\Q_\ell\to\GL_{V_\ell}$ be the algebraic morphism
associated with $\rho$.  We have
\[
\rho\circ i(x) = r(x^{-1}) \quad\text{for}\quad
x\in K_\ell^*\cap U_m = \prod_{v|\ell}U_{v,m}
\]
Define a map $\psi\colon I\to\Aut(V_\ell)$ by
\[
\psi(x) = \rho\circ i(x)\cdot r(x_\ell)
\]
where $x_\ell$ is the $\ell$-th component of the idèle $x$.  One checks
immediately that $\psi$ is \emph{trivial on} $U_m$ and \emph{coincides
with} $r$ \emph{on} $K^*$.
%% TODO-VERIFY(p105): the typescript prints a stray asterisk at subscript
%% depth after this $U_m$ (``U_{m*}''), in both the 1968 and 1989 prints;
%% no such notation is defined anywhere in the book, and the universal
%% property of $S_m$ (chap.~II, 2.2, condition (a)) needs triviality on
%% $U_m$ itself, so the mark is treated as a typist's artifact.
Hence $r$ is trivial on $E_m = K^*\cap U_m$ and factors through an
algebraic morphism $r_m\colon T_m/\Q_\ell\to\GL_{V_\ell}$.  By the
universal property of the $\Q_\ell$-algebraic group $S_m/\Q_\ell$
(cf.~chap.~\ref{chap:II}, \S\ref{sec:II.1.3} and \S\ref{sec:II.2.2}),
there exists an algebraic morphism
\[
\phi\colon S_m/\Q_\ell \to \GL_{V_\ell}
\]
with the following properties:
\begin{alist}
\item The morphism
$T_m/\Q_\ell \to S_m/\Q_\ell \xrightarrow{\ \phi\ }\GL_{V_\ell}$
is $r_m$.
\item the map
$I \xrightarrow{\ \varepsilon\ } S_m(\Q_\ell)
\xrightarrow{\ \phi\ }\Aut(V_\ell)$ is $\psi$.
\end{alist}
It is trivial to check that the \textit{l}-adic representation
$\phi_\ell$ attached to $\phi$ as above coincides with $\rho$.  Indeed,
if $a\in I$, we have (with the notations of chap.~\ref{chap:II})
\begin{align*}
\phi_\ell\circ i(a) &= \phi(\varepsilon_\ell(a))
= \phi(\varepsilon(a))\,\phi(\pi_\ell(a_\ell^{-1}))\\
&= \psi(a)\,\phi(\pi_\ell(a_\ell^{-1}))\\
&= \rho\circ i(a)\,r(a_\ell)\,\phi(\pi_\ell(a_\ell^{-1}))\\
&= \rho\circ i(a)
\end{align*}
since $\phi\circ\pi_\ell = r$ by (a) above.  Hence $\phi_\ell=\rho$; the
fact that $\rho$ is rational then implies that $\phi$ \emph{can be
defined over} $\Q$ (chap.~\ref{chap:II}, \S\ref{sec:II.2.4}, Prop.), and
this gives $V_0$ and $\phi_0$, q.e.d.
\end{proof}

\runhead{Remark}

The subspace $V_0$ of $V_\ell$ constructed in the proof of the theorem is
\emph{not} unique; however, any other choice gives us a space of the form
$\sigma V_0$, where $\sigma$ is an automorphism of $V_\ell$ commuting
with $\rho$.  To a given $V_0$ corresponds of course a unique $\phi$.

\begin{corollary}[1]\label{cor:III.2.3.1}
For each prime number $\ell'$ there exists a unique (up to isomorphism)
$\ell'$-adic rational semi-simple representation $\rho_{\ell'}$ of $K$,
compatible with $\rho$.  It is abelian and locally algebraic.  These
representations form a strictly compatible system
(cf.~chap.~\ref{chap:I}, \S\ref{sec:I.2.3}) with exceptional set contained
in $\Supp(m)$.  For an infinite number of $\ell'$, $\rho_{\ell'}$ can be
brought in diagonal form.
\end{corollary}

\begin{proof}
The unicity of the $\rho_{\ell'}$ follows from the theorem of
chap.~\ref{chap:I}, \S\ref{sec:I.2.3}.  For the existence, take
$\rho_{\ell'}$ to be the $\phi_{\ell'}$ associated to $\phi$ as in
chapter~\ref{chap:II}, \S\ref{sec:II.2.5}.  The remaining assertion
follows from the proposition in chap.~\ref{chap:II},
\S\ref{sec:II.2.5}.
\end{proof}

\begin{corollary}[2]\label{cor:III.2.3.2}
The eigenvalues of the Frobenius elements $F_{v,\rho}$
($v\notin\Supp(m)$, $p_v\neq\ell$) generate a finite extension of $\Q$.
\end{corollary}

This follows from the corresponding property of $\phi_\ell$, cf.~chapter
\ref{chap:II}, \S\ref{sec:II.2.5}, Remark~1.

\subsechead{2.4}{A mild generalization}

Most results of this and the previous Chapter may be extended to the case
where we take for ground field of the linear representation a number
field $E$ (instead of $\Q$).  More precisely, let $\lambda$ be a finite
place of $E$ and let $E_\lambda$ be the $\lambda$-adic completion of $E$.
The notion of an $E$-rational $\lambda$-adic representation of $K$ has
been defined in chap.~\ref{chap:I}, \S\ref{sec:I.2.3}, Remark.  Let
\[
\rho\colon \Gal(\bar K/K) \to \Aut(V_\lambda)
\]
be such a representation, and assume $\rho$ is abelian.  Let $\ell$ be
the residue characteristic of $\lambda$, so that $E_\lambda$ contains
$\Q_\ell$.  As in \S\ref{sec:III.2.1}, we say that $\rho$ is \emph{locally
algebraic} if there exists an algebraic morphism
\[
f\colon T/E_\lambda \to \GL_{V_\lambda}
\]
such that $\rho\circ i_\ell(x) = f(x^{-1})$ for $x\in K_\ell^*$ close
enough to $1$ (note that $K_\ell^* = T(\Q_\ell)$ is a subgroup of
$T(E_\lambda)$)
%% the original prints no period after ``T(E_\lambda))''
As in \S\ref{sec:III.2.3}, one proves that such a $\rho$ \emph{comes from
an} $E$-\emph{linear representation of some} $S_m$ (and conversely).

\subsechead{2.5}{The function field case}

The above constructions have a (rather elementary) analogue for
\emph{function fields of one variable over a finite field}:

Let $K$ be such a field, and let $p$ be its characteristic.  If $m$ is a
modulus for $K$ (i.e. a positive divisor) we define the subgroup $U_m$ of
the idèle group $I$ as in chap.~\ref{chap:II}, \S\ref{sec:II.2.1}, and we
put
\[
\Gamma_m = I/U_mK^*\,.
\]

The degree map $\mathrm{deg}\colon I\to\Z$ is trivial on $U_m$, hence
defines an exact sequence
\[
1 \to J_m \to \Gamma_m \to \Z \to 1\,.
\]

One sees easily that the group $J_m$ is \emph{finite}; moreover, it may
be interpreted as the group of rational points of the ``generalized
Jacobian variety defined by $m$''.  If $\hat{\Gamma}_m$ denotes the
completion of $\Gamma_m$ with respect to the topology of subgroups of
finite index, it is known (class field theory) that
$\Gal(\bar K/K)^{\mathrm{ab}} \simeq \varprojlim\ \hat{\Gamma}_m$.

Let now $\rho\colon \Gal(\bar K/K)^{\mathrm{ab}}\to\Aut(V_\ell)$ be an
abelian \textit{l}-adic representation of $K$, with $\ell\neq p$.  One
proves as in \S\ref{sec:III.2.2} that there exists a modulus $m$ such
that $\rho$ is trivial on $U_m$, i.e. such that $\rho$ may be identified
with a \emph{homomorphism of} $\hat{\Gamma}_m$ \emph{into}
$\Aut(V_\ell)$.  Moreover

\begin{proposition}
A homomorphism $\phi\colon \Gamma_m \to \Aut(V_\ell)$ can be extended to
a continuous homomorphism of $\hat{\Gamma}_m$ if and only if there exists
a lattice of $V_\ell$ which is stable by $\rho(\Gamma_m)$.
\end{proposition}

The necessity follows from Remark~1 of chap.~\ref{chap:I},
\S\ref{sec:I.1.1}.  The sufficiency is clear.

Note that, as in the number field case, we have Frobenius elements and we
can define the notion of \emph{rationality} of an \textit{l}-adic
representation.

\begin{theorem}
An abelian \textit{l}-adic representation
\[
\phi\colon \hat{\Gamma}_m \to \Aut(V_\ell)
\]
of $K$ is rational if and only if $\mathrm{Tr}\,\phi(\gamma)$ belongs to
$\Q$ for every $\gamma\in\Gamma_m$.
\end{theorem}

If $v\notin\Supp(m)$, and if $f_v$ is a uniformizing parameter at $v$,
the image $F_v$ of $f_v$ in $\Gamma_m$ is the Frobenius element of the
Galois group $\hat{\Gamma}_m$.  Hence, if $\mathrm{Tr}\,\phi$ takes
rational values on $\Gamma_m$, the characteristic polynomial of
$\phi(F_v)$ has rational coefficients for all $v\notin\Supp(m)$ and
$\phi$ is rational.

To prove the converse, note first that \v{C}ebotarev's theorem
(Chap.~\ref{chap:I}, \S\ref{sec:I.2.2}) is valid for $K$, if one uses a
somewhat weaker definition of equipartition.  Hence, the Frobenius
elements $F_v$ are \emph{dense} in $\hat{\Gamma}_m$.  In particular, they
generate $\Gamma_m$, and, from this, one sees that
$\mathrm{Tr}\,\rho(\gamma)$ belongs to some number field $E$.  We can
then construct an $E$-linear representation
$\phi\colon \Gamma_m\to\GL(n,E)$ with the same trace as $\rho$, and the
theorem follows from:

\begin{lemma}
Let $\Gamma$ be a finitely generated abelian group, and
$\phi\colon \Gamma\to\GL(n,E)$ a linear representation of $\Gamma$ over a
number field $E$.  Let $\Sigma$ be a subset of $\Gamma$, which is dense
in $\Gamma$ for the topology of subgroups of finite index.  Assume that
$\mathrm{Tr}\,\phi(\gamma)\in\Q$ for all $\gamma\in\Sigma$.  Then
$\mathrm{Tr}\,\phi(\gamma)\in\Q$ for all $\gamma\in\Gamma$.
\end{lemma}

\begin{proof}[Proof of the lemma.]
Since $\phi(\Gamma)$ is finitely generated, there is a finite $S$ of
places of $E$ such that all the elements of $\phi(\Gamma)$ are
$S$-integral matrices.  If $\ell'$ is a prime number not divisible by any
element of $S$, the image of $\phi(\Gamma)$ in
$\GL(n,\,E\tens\Q_{\ell'})$ is contained in a compact subgroup of
$\GL(n,\,E\tens\Q_{\ell'})$; hence $\phi$ extends by continuity to
\[
\hat{\phi}\colon \hat{\Gamma} \to \GL(n;\,E\tens\Q_{\ell'})
\]
where $\hat{\Gamma}$ is the completion of $\Gamma$ for the topology of
subgroups of finite index.  Since $\Sigma$ is dense in $\hat{\Gamma}$, it
follows that $\mathrm{Tr}\,\hat{\phi}(\hat{\gamma})$ belongs to the
adherence $\Q_{\ell'}$ of $\Q$ in $E\tens\Q_{\ell'}$ for every
$\hat{\gamma}\in\hat{\Gamma}$.  Hence, if $\gamma\in\Gamma$, we have
\[
\mathrm{Tr}\,\phi(\Gamma) \in E\cap\Q_{\ell'} = \Q\,,
\qquad\text{q.e.d.}
\]
%% the display argument is printed ``Tr phi(Gamma)'' (capital Gamma) in
%% the original, although a single ``gamma'' is meant.
\end{proof}

\begin{exercises}
\begin{nlist}
\item Let $\phi\colon \Gamma_m\to\Aut(V_\ell)$ be a semi-simple
\textit{l}-adic representation of $\Gamma_m$.  Show the equivalence of:
\begin{alist}
\item $\phi$ extends continuously to $\hat{\Gamma}_m$.
\item For every $\gamma\in\Gamma_m$, the eigenvalues of $\phi(\gamma)$
are units (in a suitable extension of $\Q_\ell$).
\item There exists $\gamma\in\Gamma_m$, with $\mathrm{deg}(\gamma)\neq 0$,
such that the eigenvalues of $\phi(\gamma)$ are units.
\item For every $\gamma\in\Gamma_m$, one has
$\mathrm{Tr}\,\phi(\gamma)\in\Z_\ell$.
\end{alist}
\item Let $\phi\colon \hat{\Gamma}_m\to\Aut(V_\ell)$ be a rational
\textit{l}-adic representation of $K$.  Show that, for almost all prime
number $\ell'$, there is a rational $\ell'$-adic representation of $K$
compatible with $\phi$.  Show that this holds for all $\ell'\neq p$ if
and only if the following property is valid: for all
$\gamma\in\Gamma_m$, the coefficients of the characteristic polynomial of
$\phi(\gamma)$ are $p$-integers.
\end{nlist}
\end{exercises}
